% !PRINTED 532
\weberpage{532}

From (27) one derives a simpler formula for the function
\[
f(\omega) = e^{-\frac{\pi i}{24}}\,\frac{\eta\!\left(\frac{\omega+1}{2}\right)}{\eta(\omega)} = e^{\frac{\pi i}{24}}\,\frac{\eta\!\left(\frac{\omega-1}{2}\right)}{\eta(\omega)}.
\]

For if one replaces
\[
\begin{aligned}
A &\ \text{by}\ 2A,\\
B &\ \text{\textquotedblbase}\ \ A+B,\\
C &\ \text{\textquotedblbase}\ \ \tfrac12(A+2B+C),
\end{aligned}
\]
then $m=AC-B^2$ remains unchanged, and $\omega_1,\omega_2$ pass over into
\[
\omega_1' = \frac{\omega_1+1}{2},\quad \omega_2' = \frac{\omega_2-1}{2}.
\]

If one therefore sets
\weq{29}{\begin{aligned}
\psi &= (A,\,2B,\,C),\\
\psi' &= [2A,\,2(A+B),\,\tfrac12(A+2B+C)],
\end{aligned}}
then there results from (27)
\weq{30}{\mathop{\mathrm{Lim}}\limits_{s=1}\left(\sum^{x,y}\frac{1}{\psi^s} - \sum^{x,y}\frac{1}{\psi'^s}\right) = \frac{2\pi}{\sqrt m}\log\frac{f(\omega_1)f(\omega_2)}{\sqrt 2}.}

We now take not only $A, 2B, C$, but also $A, B, C$, to be integers without common divisor, and $A$ odd. Then $\psi$ is a primitive form of discriminant $-4m$, and we consider first the two cases
\[
m\equiv 1,\quad m\equiv 2\ (\mathrm{mod}\ 4).
\]

If we assume, which is no further restriction, that $B\equiv m+1\ (\mathrm{mod}\ 2)$, that is, $B$ even or odd according as $m$ is odd or even, then there results from
\[
m = AC-B^2\text{:}
\]
\[
\begin{aligned}
A-C&\equiv 0\ (\mathrm{mod}\ 4) \quad m\ \text{odd},\\
A+C&\equiv 0\ \text{\textquotedblbase} \qquad\ \ m\ \text{even},
\end{aligned}
\]
and the form $\psi'$ is primitive of discriminant $-4m$. By the composition of quadratic forms we have:
\[
\psi' = \psi_0\psi,
\]
if
\[
\begin{aligned}
\psi_0 &= \left(2,\,2,\,\tfrac{m+1}{2}\right) & (m\ \text{odd}),\\
&= \left(2,\,0,\,\tfrac{m}{2}\right) & (m\ \text{even}),
\end{aligned}
\]
%CONT>>
